\documentclass{article}
\usepackage[utf8]{inputenc}
\usepackage[T1]{fontenc}
\usepackage{microtype}
\usepackage{graphicx}
\usepackage{lipsum}
\usepackage{newspaper}
%that's the one which makes the magic happen
\usepackage{newspaper-mod}
%%to redefine the headline and byline style if you want to.
%% These can be issued just before any
%% byline or headline in the paper, to
%% individually style each article
%%
% \renewcommand{\headlinestyle}{\itshape\Large\lsstyle}
% \renewcommand{\bylinestyle}{\bfseries\Large\raggedright}

\usepackage{multicol}
%For making columns using command \begin{multicols}{2}....\end{multicols}

%Fancy images
\usepackage{picinpar}
%uasage of picinpar:
%\begin{window}[1,l,\includegraphics{},caption]xxxxx\end{window}

% \begin{window}[2,r,\includegraphics[width=1.0in]{atom.jpg},\centerline{The Atom}] The \verb+multicol+ package allows using multiple columns without starting a new page.  Using floats is not possible in a columns environment, however with the \verb+picinpar+ package, I can set a picture inside a block of text---just like you one you see here.  Isn't \LaTeX{} cool?
% And now we're just filling more space, and yet more space.  
% \end{window}


% TODO: Change the font to Times New Roman



%Metadata you shouldn't change
\SetPaperName{
          %  \includegraphics[scale=0.3]{fps}
        % \scalerel*{\includegraphics{fps}}{f}
            FPS Eagles:}
\SetHeaderName{FPS Eagles}
\SetPaperLocation{OLDC Karachi}
\SetPaperSlogan{
            %     \includegraphics[scale=0.3]{fps}
                ``\textit{FPSonians are the best; second to none!}''}
\SetPaperPrice{Founded by Awab Qureshi and Sukaina Kazmi; batch 2020}



% Metadata you should change
\date{\today}
\currentvolume{1}
\currentissue{1}




\begin{document}
\maketitle
\large 

%The beginning 

%                               FRONT PAGE

\headline{PLASTIC}\textbf{by Awab Ahmed Qureshi XI-D}

Researchers recently probed fish-life in the deepest parts of our ocean and found something shocking. In their digestive tracts were quantities of microplastics. The study raised several concerns. Finding microplastics at such a depth means that even the deepest parts of our oceans are not safe from our litter.More and more plastic is thrown into our oceans every day with some estimates ranging between 4.8~12.7 million tonnes per Annum.

Indeed microplastics are quite dangerous as the ones found in such fish could easily transfer up the food chain and into our food supply, causing a plethora of adverse health effects.

Students in the SLC in FPS realised this crisis, and thus, to prevent it, worked with the admin to ban the sale of all plastic bottles in the canteen. The ban, which was put in place earlier this term has reduced the plastic bottle waste OLDC generates to $\boldsymbol{0}$. In place of this, students can now refill their bottles at a lower rate from the canteen's water filter.

It is hoped that our contribution, whilst small, will make a difference by discouraging the use of disposable plastics and will a step on the long road towards a green future. 


\closearticle
%do \closearticle for that double line bar

\begin{multicols}{2}
%to make columns

\headline{Escape to Wonderland}\textbf{Sukaina Kazmi XI-C}

Friday 11th October was a night to remember for all grade 11s. FPS A-Levels had held its 4th annual “Escape to Wonderland” which unfortunately the last of its kind. The activity aided in enhancing and bringing out a student’s team building and problem solving skills as the students hurriedly pieced together multiple clues and puzzles to save the world from the apocalypse .Students in teams of eight were locked in a military bunker with an attendant and had approximately an hour to stop the virus from spreading into the atmosphere and escaping the room. Students had to channel their cognitive skills to scavenge the keys to the buzzer thereby emerging triumphant. Once the students succeeded in escaping the room they gained entry to the magical wonderland full of food, games, rides and music for them to dance the night away. The night ended up being a truly memorable event which I’m sure this batch is to remember.


\headline{FPS Ignite}\textbf{Laiba Abid XI-C}

Ignite is an annual three day event hosted by FPS A Levels and provides students a platform to test out their entrepreneurial ideas. Participants form 4-member delegations and the shortlisted ones then partake in bidding for their desired industry. Upon receiving the bidding results, they then design a pitch presentation, and legal documents for their product. 

On the first day delegations were tasked with designing billboards. On the second day they raced to solve a crisis and shoot their TV commercial. In the evening, participants got to socialize with their batch-mates and blow off some steam in the Social. The third day kept everyone on their toes as the delegations had to go through a ‘surprise round’, pitch their product to the investors and present their tvc and billboard. The event ended with a grand award ceremony where various trophies were given out.

\end{multicols}
\pagebreak 
%to make sure stuff stays on this page


%                               2ND PAGE
\headline{Sports}
\begin{multicols}{2}

Sarah Sajjad, our very own pride of FPS, as well as an ambitious champion to represent the province of Sindh in the National Olympic Games held in Peshawar, and impressing us once again as the youngest throwball athlete of the event. However, this is not the first time Sara has made FPS and us FPSonians proud; infact, this outstanding athletic virtuoso has made our school appreciative several times last year in inter-school competitions. She left for Peshawar, along with her mother for the nerve wrecking journey, on Nov, 8th. The team started their first match with Balochistan, winning it with a lead of 20 points. However her last match made her quite nervous, she knew she didn’t have to let it affect her game and strength. Among the 6 games that were held, the two qualifying matches, two of the semi-finals and two of the finals, Sindh won all 6 of them. Sara honored us all, as well as her family, and most importantly her parents, with her trophy and gold medal.  
\end{multicols}
\closearticle
\pagebreak

%                              3RD PAGE 

\begin{multicols}{2}

\headline{This Week in History}
1912-The ocean liner Titanic sinks after hitting an iceberg in the North Atlantic Sea.

1966-Football’s World Cup is stolen from a display in London. A dog finds it a week later, under a hedge.

1968-The funeral of Martin Luther King Jr. takes place.

1973-The art world mourns Spanish artist Pablo Picasso who dies in France aged 91.
\closearticle

\headline{Upcoming Events}
THE WINTER FEST ~ 19th December 2019

\pagebreak


\end{multicols}

\end{document}